\chapter{Image-Guided Biopsy}

\label{ch:ch77_biopsy}

Image-guided biopsy is a minimally invasive procedure in which a needle is directed into a lesion under real-time imaging (CT, ultrasound, or fluoroscopy) to obtain tissue for diagnosis. It avoids open surgery in many cases and provides definitive histopathological diagnosis to guide treatment.

---

\section{1. Patient Perspective: What a Needle Biopsy Involves}

If your doctor needs to know exactly what a lump or abnormal area is made of, you may be offered a \textbf{needle biopsy}. This means using a thin needle to take a small sample of tissue.

\begin{itemize}
  \item The procedure is performed under local anaesthetic (injection to numb the skin and deeper tissues).
  \item An imaging scan (CT or ultrasound) is used to guide the needle precisely into the target.
  \item You may feel pressure or mild discomfort but should not feel sharp pain.
  \item For lung biopsies, you will be asked to hold your breath briefly. For liver or kidney biopsies, you may be asked to breathe normally or hold your breath at specific moments.
  \item The procedure typically takes 20--45 minutes.
  \item After the biopsy, you will be observed for 2--4 hours (lung biopsies may require overnight observation) to watch for complications such as bleeding or air leak.
\end{itemize}

---

\section{2. Indications and Pre-Procedure Planning}

\begin{itemize}
  \item \textbf{Indications}: Undiagnosed mass or lesion, staging of known malignancy, suspected infection requiring microbiological diagnosis, monitoring of treatment response, or assessment of transplant rejection.
  \item \textbf{Contraindications}: Uncorrectable coagulopathy (INR $>$ 1.5, platelets $<$ 50,000), unsafe access route, uncooperative patient, suspected vascular lesion (e.g., pheochromocytoma without preparation), and pneumothorax risk in apical lung lesions without surgical backup.
\end{itemize}

\subsection{Route Planning}
\begin{itemize}
  \item Choose the \textbf{shortest, safest route} that avoids major vessels, nerves, bowel, and the pleura (for lung biopsies).
  \item For lung biopsies, a \textbf{transthoracic approach} is used. The needle traverses the pleura; the risk of pneumothorax must be weighed against the diagnostic need.
  \item For deep abdominal lesions, avoid transgressing the kidney, bowel, or gallbladder.
  \item The route should \textbf{preserve future treatment options} (e.g., avoid seeding a biopsy tract in a resectable hepatocellular carcinoma if possible).
\end{itemize}

\section{3. CT vs Ultrasound Guidance}

\begin{table}[h]
\centering
\begin{tabular}{l l l}
\hline
\textbf{Feature} & \textbf{CT Guidance} & \textbf{Ultrasound Guidance} \\
\hline
Accessibility & Requires CT suite & Portable, bedside \\
Best for & Deep lesions, lung, bone & Superficial lesions, liver, thyroid, breast \\
Cost & Higher & Lower \\
Radiation & Yes (proportional to scans) & None \\
Real-time & Intermittent (scan, advance, scan) & Continuous real-time \\
\hline
\end{tabular}
\caption{CT vs Ultrasound Guidance for Biopsy}
\end{table}

\section{4. Core Needle vs Fine Needle Aspiration}

\begin{itemize}
  \item \textbf{Fine needle aspiration} (FNA, 20--25G): Obtains cells for cytology. Rapid, lower complication rate, but may not provide sufficient tissue for immunohistochemistry or molecular studies.
  \item \textbf{Core needle biopsy} (CNB, 14--18G): Obtains a tissue core for histology. Provides architectural information and more tissue for ancillary testing. Higher diagnostic yield for most solid organ lesions.
  \item \textbf{Choice depends on}: Lesion type, organ, clinical question, and available pathology expertise. Core biopsy is generally preferred for solid organ masses; FNA is often sufficient for lymph nodes, thyroid nodules, and suspected metastatic disease.
\end{itemize}

\begin{pearlbox}[Specimen Handling]
Label all specimens immediately with the correct patient details and site. Place core biopsies in formalin and FNA specimens in saline or cytology fixative. If ROSE (rapid on-site evaluation) is available, the cytopathologist can assess adequacy during the procedure, reducing the need for repeat biopsies.
\end{pearlbox}

---

\section{5. Complications}

\begin{table}[h]
\centering
\begin{tabular}{l l l}
\hline
\textbf{Complication} & \textbf{Organ at Risk} & \textbf{Management} \\
\hline
Haemorrhage & Any solid organ & Embolisation or surgical backup \\
Infection & Any site & Antibiotics \\
Needle tract seeding & Hepatocellular carcinoma & Minimise passes; consider RFA tract ablation \\
Pneumothorax & Lung biopsy & Observation or chest tube \\
Vasovagal reaction & Any site & Positional support, atropine if needed \\
\hline
\end{tabular}
\caption{Complications of Image-Guided Biopsy}
\end{table}

\subsection{When to Stop: The CT-Guided Lung Biopsy}
During a CT-guided lung biopsy, a post-biopsy scan is performed. If a significant pneumothorax develops (lung $>$ 25\% collapsed, or the patient is symptomatic), the procedure is stopped and the patient is monitored. Small, asymptomatic pneumothoraces can be observed; larger ones may require aspiration or chest tube insertion.

\section{6. Sample Procedure Note}

\begin{testbox}[Procedure: CT-Guided Percutaneous Core Needle Biopsy]
\textbf{Date}: [Date] \quad \textbf{Operator}: [Name, grade] \quad \textbf{Assistant}: [Name]\\
\textbf{Consent}: Informed consent obtained. Risks of bleeding, pneumothorax, infection, and damage to adjacent structures discussed.\\
\textbf{Indication}: 3.5 cm right lower lobe mass. Staging CT shows no metastases. Tissue diagnosis required.\\
\textbf{Anaesthesia}: Local anaesthetic (1\% lidocaine 10 mL) and conscious sedation (midazolam 2 mg, fentanyl 50 mcg IV).\\
\textbf{Technique}: Patient positioned prone on CT table. CT performed to confirm lesion location. Skin prepared and draped. Local anaesthetic infiltrated. 18G co-axial core biopsy needle advanced under intermittent CT guidance through the posterior chest wall into the lesion. Three core samples obtained. Post-biopsy CT shows no pneumothorax.\\
\textbf{Specimens}: Three core biopsies placed in 10\% neutral buffered formalin. Sent to histopathology.\\
\textbf{Complications}: None.\\
\textbf{Post-procedure}: Patient transferred to recovery. Observed for 4 hours with serial CXR. Discharged if stable.
\end{testbox}

\clearpage
